\documentclass[10pt,a4paper]{amsart}
\usepackage{amsmath,amssymb,amsthm,graphicx}
\usepackage{verbatim}
\newenvironment{reference}{\comment}{\endcomment}
%\newenvironment{reference}{}{}
\newenvironment{slogan}{\comment}{\endcomment}
\newenvironment{history}{\comment}{\endcomment}

% For commutative diagrams we use Xy-pic
\usepackage[all]{xy}

% We use 2cell for 2-commutative diagrams.
\xyoption{2cell}
% Theorem environments.
%
\theoremstyle{plain}
\newtheorem{theorem}[subsection]{Theorem}
\newtheorem{proposition}[subsection]{Proposition}
\newtheorem{lemma}[subsection]{Lemma}

\theoremstyle{definition}
\newtheorem{definition}[subsection]{Definition}
\newtheorem{example}[subsection]{Example}
\newtheorem{exercise}[subsection]{Exercise}
\newtheorem{situation}[subsection]{Situation}

\theoremstyle{remark}
\newtheorem{remark}[subsection]{Remark}
\newtheorem{remarks}[subsection]{Remarks}

\numberwithin{equation}{subsection}

% Macros
%
\def\lim{\mathop{\mathrm{lim}}\nolimits}
\def\colim{\mathop{\mathrm{colim}}\nolimits}
\def\Spec{\mathop{\mathrm{Spec}}}
\def\Hom{\mathop{\mathrm{Hom}}\nolimits}
\def\Ext{\mathop{\mathrm{Ext}}\nolimits}
\def\SheafHom{\mathop{\mathcal{H}\!\mathit{om}}\nolimits}
\def\SheafExt{\mathop{\mathcal{E}\!\mathit{xt}}\nolimits}
\def\Sch{\mathit{Sch}}
\def\Mor{\mathop{\mathrm{Mor}}\nolimits}
\def\Ob{\mathop{\mathrm{Ob}}\nolimits}
\def\Sh{\mathop{\mathit{Sh}}\nolimits}
\def\NL{\mathop{N\!L}\nolimits}
\def\CH{\mathop{\mathrm{CH}}\nolimits}
\def\proetale{{pro\text{-}\acute{e}tale}}
\def\etale{{\acute{e}tale}}
\def\QCoh{\mathit{QCoh}}
\def\Ker{\mathop{\mathrm{Ker}}}
\def\Im{\mathop{\mathrm{Im}}}
\def\Coker{\mathop{\mathrm{Coker}}}
\def\Coim{\mathop{\mathrm{Coim}}}

% Boxtimes
%
\DeclareMathSymbol{\boxtimes}{\mathbin}{AMSa}{"02}

%
% Macros for moduli stacks/spaces
%
\def\QCohstack{\mathcal{QC}\!\mathit{oh}}
\def\Cohstack{\mathcal{C}\!\mathit{oh}}
\def\Spacesstack{\mathcal{S}\!\mathit{paces}}
\def\Quotfunctor{\mathrm{Quot}}
\def\Hilbfunctor{\mathrm{Hilb}}
\def\Curvesstack{\mathcal{C}\!\mathit{urves}}
\def\Polarizedstack{\mathcal{P}\!\mathit{olarized}}
\def\Complexesstack{\mathcal{C}\!\mathit{omplexes}}
% \Pic is the operator that assigns to X its picard group, usage \Pic(X)
% \Picardstack_{X/B} denotes the Picard stack of X over B
% \Picardfunctor_{X/B} denotes the Picard functor of X over B
\def\Pic{\mathop{\mathrm{Pic}}\nolimits}
\def\Picardstack{\mathcal{P}\!\mathit{ic}}
\def\Picardfunctor{\mathrm{Pic}}
\def\Deformationcategory{\mathcal{D}\!\mathit{ef}}
\usepackage{cleveref}
\crefname{originallemma}{original lemma}{original lemmas}
\Crefname{originallemma}{Original lemma}{Original lemmas}
\makeatletter
\newlabel{lemma-flat-local-given-residue-field}{{03C3}{source}{}{sourcepointer-1}{}}
\newlabel{lemma-flat-local-given-residue-field@cref}{{[originallemma][0][]03C3}{[1][0][]source}{}{}{}}
\makeatother
\title{A separable algebraic residue field extension is realised by a colimit of finite etale extensions}
\author{The Stacks Project Authors}
\date{stacks-project:algebra.tex:local-snapshot-20261012; GNU FDL 1.2 or later}
\begin{document}
\maketitle
\noindent\textit{Editorial note.} All delivered mathematical text is the author's own, copied from the pinned Stacks Project chapter file algebra.tex; no statement or proof step was rewritten, reordered or removed.\par\medskip
\noindent\textit{Editorial note.} Preamble provenance, as required for this source: the frozen snapshot stages the chapter's own preamble file beside it (preamble.tex, 217 lines, 8116 bytes, the same shared Stacks-wide preamble the published chow and etale-cohomology cards use); the first line of the pinned chapter file is \textbackslash{}input\{preamble\}, so no preamble text lives inside the chapter itself. The card takes that file's lines 8-18 (the verbatim package and the three \textbackslash{}comment-based environments reference/slogan/history at lines 9-12, section 13.2, together with the Xy-pic setup) and its lines 154 to the runtime end of the file (the native theorem declarations and mathematical macros). Not one character was changed, no package was added, no macro was defined, and the excerpt serves only as preamble, never as a span coordinate.\par\medskip
\noindent\textit{Editorial note.} The same-chapter pointers are located at their \textbackslash{}label lines inside the pinned chapter file; the cross-chapter pointers are located in the project's own tag list staged in the snapshot (sources/stacks-chow/tags.txt), at the verbatim lines carrying their canonical tags.\par\medskip
\noindent\textit{Editorial note.} The counted proof refers to the earlier construction by label; it is recorded as a navigation stub pointed at its \textbackslash{}label line, and the argument consuming it stays inside the counted span.\par\medskip
\noindent\textit{Editorial note.} The counted spans contain no \textbackslash{}cite, so the card carries no bibliography.\par\medskip
\noindent\textit{Allowed standard prerequisites (editorial).}\begin{itemize}
\item The extension constructed for a prescribed residue field. Used at the start of the counted proof (lines 44869-44874). Provides a well-ordered colimit presentation whose successor maps are finite etale.
\item Transfinite induction over the well-ordered index set. Used throughout the counted proof (lines 44874-44888). Reduces the statement to successor and limit steps, each handled by a separate argument.
\end{itemize}\medskip
\ifdefined\hypertarget\hypertarget{sourcepointer-1}{}\fi
\noindent\textit{Source reference.} Original lemma 03C3: sources/stacks-algebra/pinned-original-source.tex, source lines 44790--44790. Navigation pointer for \textbackslash{}ref\{lemma-flat-local-given-residue-field\} used inside the counted span (line 44870); the argument that consumes it is retained in the span.\par\medskip
\begin{lemma}
\label{lemma-colimit-finite-etale-given-residue-field}
Let $(R, \mathfrak m, k)$ be a local ring. If $k \subset K$ is a
separable algebraic extension, then there exists a directed set $I$ and
a system of finite \'etale extensions $R \subset R_i$, $i \in I$
of local rings such that $R' = \colim R_i$ has residue field
$K$ (as extension of $k$).
\end{lemma}

\begin{proof}
Let $R \subset R'$ be the extension constructed in the proof of
Lemma \ref{lemma-flat-local-given-residue-field}. By construction
$R' = \colim_{\alpha \in A} R_\alpha$ where $A$ is a well-ordered
set and the transition maps $R_\alpha \to R_{\alpha + 1}$
are finite \'etale and $R_\alpha = \colim_{\beta < \alpha} R_\beta$
if $\alpha$ is not a successor. We will prove the result by transfinite
induction.

\medskip\noindent
Suppose the result holds for $R_\alpha$, i.e., $R_\alpha = \colim R_i$
with $R_i$ finite \'etale over $R$. Since
$R_\alpha \to R_{\alpha + 1}$ is finite \'etale
there exists an $i$ and a finite \'etale extension $R_i \to R_{i, 1}$
such that $R_{\alpha + 1} = R_\alpha \otimes_{R_i} R_{i, 1}$.
Thus $R_{\alpha + 1} = \colim_{i' \geq i} R_{i'} \otimes_{R_i} R_{i, 1}$
and the result holds for $\alpha + 1$. Suppose $\alpha$ is not a successor
and the result holds for $R_\beta$ for all $\beta < \alpha$.
Since every finite subset $E \subset R_\alpha$ is contained in $R_\beta$
for some $\beta < \alpha$ and we see that $E$ is contained in a finite \'etale
subextension by assumption. Thus the result holds for $R_\alpha$.
\end{proof}

\end{document}
